% Problem-section file following ../_template.tex.
\section{Candidate pure-loss second-order converse}
\label{sec:problem-76}

\paragraph{Problem.}
Does the pure-loss bosonic channel admit the following candidate second-order
classical converse under a maximum-photon-number occupation constraint?  Let
$\mathcal N_\eta$ be the single-mode pure-loss channel with transmissivity
$0<\eta<1$, defined in the Heisenberg picture by
\begin{equation}
  \hat b=\sqrt{\eta}\,\hat a+\sqrt{1-\eta}\,\hat e,
  \label{eq:p76-pure-loss-channel}
\end{equation}
where the environment mode $E$ is in the vacuum.  In an $n$-use code, the
channel in Eq.~\eqref{eq:p76-pure-loss-channel} is used to transmit one of
$M$ input states $\rho_m^{A^n}$, with a decoding POVM
$\{\Lambda_m^{B^n}\}_{m=1}^M$.  Write
$\overline\rho_{A^n}:=M^{-1}\sum_m\rho_m^{A^n}$ and let
$\Pi_{\lceil nN_S\rceil}$ project onto the $n$-mode subspace of total photon
number at most $\lceil nN_S\rceil$.  For fixed $N_S>0$,
$\varepsilon\in(0,1)$, and $c>0$, impose
\begin{equation}
  \frac1M\sum_{m=1}^M
    \operatorname{Tr}\!\left[
      \Lambda_m\mathcal N_\eta^{\otimes n}(\rho_m)
    \right]
    \geq1-\varepsilon,
  \qquad
  \operatorname{Tr}\!\left[
    \Pi_{\lceil nN_S\rceil}\overline\rho_{A^n}
  \right]
    \geq1-\delta_n,
  \qquad
  0\leq\delta_n\leq2^{-cn}.
  \label{eq:p76-code-constraints}
\end{equation}
Let $M^*_{\rm occ}(n,\eta,N_S,\varepsilon,c)$ be the largest $M$ satisfying
Eq.~\eqref{eq:p76-code-constraints}.  Define the thermal entropy and its
entropy variance by
\begin{equation}
  g(x):=(x+1)\log_2(x+1)-x\log_2x,
  \qquad
  v(x):=x(x+1)
    \left[\log_2(x+1)-\log_2x\right]^2,
  \label{eq:p76-thermal-entropy-and-variance}
\end{equation}
where $0\log_2 0:=0$.  With the functions in
Eq.~\eqref{eq:p76-thermal-entropy-and-variance}, is the following upper bound
valid as $n\to\infty$?
\begin{equation}
  \log_2 M^*_{\rm occ}(n,\eta,N_S,\varepsilon,c)
  \leq ng(\eta N_S)
    +\sqrt{n\,v(\eta N_S)}\,\Phi^{-1}(\varepsilon)
    +O(\log n),
  \label{eq:p76-second-order-converse}
\end{equation}
Here $\Phi^{-1}$ in Eq.~\eqref{eq:p76-second-order-converse} is the inverse
standard-normal cumulative distribution function, and the implicit constant
may depend on $\eta,N_S,\varepsilon,$ and $c$, but not on $n$.

\subsection*{Status}

Unsolved

\subsection*{Source}

Wilde, Renes, and Guha leave open whether their Gaussian second-order
expression is an upper bound and observe that such a converse should impose a
photon-number occupation constraint similar to Wilde and Winter's.
Equation~\eqref{eq:p76-code-constraints} adopts a particular exponentially
strong version of that suggestion, so Eq.~\eqref{eq:p76-second-order-converse}
is a precise strengthened formulation rather than a verbatim conjecture from
their paper
\sourcecite{ref:p76-wilde-renes-guha}{WRG16}.

\subsection*{Progress}

\begin{itemize}
  \item Wilde, Renes, and Guha proved the corresponding second-order
  achievability expression
  \begin{equation}
    \log_2 M^*(\mathcal N_\eta^{\otimes n},N_S,\varepsilon)
    \geq ng(\eta N_S)
      +\sqrt{n\,v(\eta N_S)}\,\Phi^{-1}(\varepsilon)
      +O(\log n).
    \label{eq:p76-second-order-achievability}
  \end{equation}
  Their basic derivation of
  Eq.~\eqref{eq:p76-second-order-achievability} uses a mean-photon-number
  ensemble.  They also modify the construction to obtain exponentially small
  leakage outside the cutoff subspace, but with degraded second-order
  parameters; neither argument supplies
  Eq.~\eqref{eq:p76-second-order-converse}
  \sourcecite{ref:p76-wilde-renes-guha}{WRG16}.

  \item Wilde and Winter proved a first-order strong converse at the rate
  $g(\eta N_S)$ under the maximum-photon-number occupation constraint.  For
  every fixed energy and rate backoff, their modified coherent-state codes
  have exponentially small cutoff leakage.  This does not establish the
  exact $g(\eta N_S)$ first-order term and the stated dispersion while also
  maintaining an arbitrary prescribed exponent $c$ in
  Eq.~\eqref{eq:p76-code-constraints}.  Moreover, a first-order strong
  converse controls rates a fixed distance above capacity and does not
  determine the
  $\sqrt n$ coefficient in Eq.~\eqref{eq:p76-second-order-converse}
  \sourcecite{ref:p76-wilde-winter}{WW14}.

  \item Under only a mean-photon-number constraint, Wilde and Winter
  constructed codes with nonvanishing success probability at rates above
  $g(\eta N_S)$.  Hence a converse of the form
  Eq.~\eqref{eq:p76-second-order-converse} is false under that weaker
  constraint; the occupation hypothesis in
  Eq.~\eqref{eq:p76-code-constraints} is essential
  \sourcecite{ref:p76-wilde-winter}{WW14}.
\end{itemize}

\subsection*{References}

\begin{enumerate}[label={},leftmargin=4.8em,labelsep=0.5em]
  \footnotesize
  \item[\textup{[WRG16]}]\label{ref:p76-wilde-renes-guha}
  M. M. Wilde, J. M. Renes, and S. Guha, ``Second-Order Coding Rates for
  Pure-Loss Bosonic Channels,'' \emph{Quantum Information Processing}
  \textbf{15}, 1289--1308 (2016).
  \href{https://doi.org/10.1007/s11128-015-0997-x}%
       {doi:10.1007/s11128-015-0997-x};
  \href{https://arxiv.org/abs/1408.5328}{arXiv:1408.5328}.

  \item[\textup{[WW14]}]\label{ref:p76-wilde-winter}
  M. M. Wilde and A. Winter, ``Strong Converse for the Classical Capacity of
  the Pure-Loss Bosonic Channel,'' \emph{Problems of Information
  Transmission} \textbf{50}, 117--132 (2014).
  \href{https://doi.org/10.1134/S003294601402001X}%
       {doi:10.1134/S003294601402001X};
  \href{https://arxiv.org/abs/1308.6732}{arXiv:1308.6732}.
\end{enumerate}

\subsection*{Comment}

The unresolved task is the candidate sharp Gaussian $\sqrt n$ converse term
under the stated fixed-exponent occupation constraint, not the first-order
capacity or strong converse.  No matching achievability theorem is known for
this exact constraint.  The exclusion of $\eta=1$ is essential: at the
identity-channel endpoint, the dimension of the cutoff Fock subspace itself
violates Eq.~\eqref{eq:p76-second-order-converse} for
$\varepsilon<\tfrac12$.

\subsection*{Tag}

Bosonic channels; Classical capacity; Finite-blocklength quantum information;
Strong converse; Continuous-variable systems; Quantum Shannon theory

\subsection*{ID}

\texttt{op\_89fb664ba06ba5de}
